% Formalization-oriented distillation of Davis--Kahan (1970), Part III.
% This is an independent mathematical writeup, not a transcription.
% Model provenance for this distillation and the associated Lean scaffold:
% GPT 5.6 High.

\documentclass[11pt]{article}
\usepackage[T1]{fontenc}
\usepackage{lmodern}
\usepackage[margin=1in]{geometry}
\usepackage{amsmath,amssymb,amsthm,mathtools,booktabs,longtable}
\usepackage[colorlinks=true,linkcolor=black,citecolor=blue,urlcolor=blue]{hyperref}

\newcommand{\K}{\mathbb K}
\newcommand{\Hh}{\mathcal H}
\newcommand{\Rr}{\mathcal R}
\newcommand{\Nn}{\mathcal N}
\newcommand{\spec}{\operatorname{spec}}
\newcommand{\ran}{\operatorname{ran}}
\newcommand{\dist}{\operatorname{dist}}
\newcommand{\diag}{\operatorname{diag}}
\newcommand{\pinch}{\operatorname{pinch}}
\newcommand{\sv}{\operatorname{sv}}
\newcommand{\id}{\operatorname{id}}
\newcommand{\op}{\mathrm{op}}
\newcommand{\HS}{\mathrm{HS}}
\newcommand{\abs}[1]{\left|#1\right|}
\newcommand{\norm}[1]{\left\lVert#1\right\rVert}
\newcommand{\inner}[2]{\left\langle#1,#2\right\rangle}

\newtheorem{theorem}{Theorem}[section]
\newtheorem{lemma}[theorem]{Lemma}
\newtheorem{proposition}[theorem]{Proposition}
\newtheorem{definition}[theorem]{Definition}
\theoremstyle{remark}
\newtheorem{remark}[theorem]{Remark}

\title{The finite-dimensional Davis--Kahan theory:\\
formalization-oriented proof distillation and Lean roadmap}
\author{GPT 5.6 High}
\date{10 July 2026}

\begin{document}
\maketitle

\begin{abstract}
This document distills the finite-dimensional content and proof architecture of
Davis and Kahan's \emph{The Rotation of Eigenvectors by a Perturbation. III}
(SIAM Journal on Numerical Analysis 7 (1970), 1--46), together with the
finite-dimensional antecedents in Davis (1963).  It is written to support the
Lean files under
\texttt{ForMathlib/Analysis/InnerProductSpace/DavisKahanTheory/}.
It deliberately expands the operator-theoretic prose into the intermediate
objects, block identities, singular-value statements, and spectral-selection
lemmas that a final Lean proof will require.  It is an independent distillation,
not a transcription of the source paper.
\end{abstract}

\tableofcontents

\section{Scope and source map}

The 1970 paper centers on four independent sharp estimates:
\(\sin\Theta\), \(\tan\Theta\), \(\sin 2\Theta\), and
\(\tan 2\Theta\).  Each has a perturbation form and a residual form, and the
bounds are stated for arbitrary unitarily invariant norms.  The paper separates
its development into: the geometry of two subspaces; the direct rotation and
its extremal properties; a sharp estimate for the Sylvester expression
\(AX-XB\); proofs of the single-angle theorems; proofs of the double-angle
theorems; and the spectral-selection/repulsion interpretation of the
double-angle bounds.

The Lean scaffold follows the same dependency order.  Definitions are made on
submodules and isometric embeddings rather than on a chosen pair of coordinate
bases.  Finite dimensionality permits every spectral statement to be reduced to
an orthonormal eigenbasis, but basis independence must be proved before the
public theorem layer is considered complete.

\section{Coordinate-free reducing subspaces and spectral blocks}

Let \(A:\Hh\to\Hh\) be Hermitian and let \(P\) be the orthogonal projector
onto a reducing subspace \(U\).  Thus
\[
  AP=PA, \qquad A=PAP+P^\perp A P^\perp.
\]
In finite dimensions it is convenient to define the spectrum carried by
\(U\) as the set of real \(\lambda\) for which there is a nonzero
\(x\in U\) with \(Ax=\lambda x\).  The canonical spectral subspace for a set
\(\Omega\subseteq\mathbb R\) is
\[
  U_\Omega(A)=\operatorname{span}\{x:Ax=\lambda x,\ \lambda\in\Omega\}.
\]
The associated orthogonal projector is denoted \(P_\Omega(A)\).

The formalization must prove that this construction is independent of the
chosen eigenbasis, reduces \(A\), has the expected range and kernel, and
respects complements, disjoint unions, intersections, and restriction to
intervals.  Repeated eigenvalues must not be split by an arbitrary ordering of
an eigenbasis.

For two blocks \((A,U)\) and \((B,V)\), the basic gap predicates are:
\begin{align*}
  \operatorname{Sep}_\delta(A|_U,B|_V)
    &:\Longleftrightarrow
      \abs{\lambda-\mu}\ge\delta
      \quad(\lambda\in\spec(A|_U),\ \mu\in\spec(B|_V)),\\
  \operatorname{OrdSep}_\delta(A|_U,B|_V)
    &:\Longleftrightarrow
      \lambda+\delta\le\mu
      \quad(\lambda\in\spec(A|_U),\ \mu\in\spec(B|_V)).
\end{align*}
The sharp single-angle theorem uses interval/exterior or ordered separation;
a merely disjoint pair of arbitrary spectral sets leads instead to the later
\(\pi/2\)-constant extension.

\section{Unitarily invariant norms, singular values, and rectangular maps}

For a linear map \(X:E\to F\) between finite-dimensional Hilbert spaces, let
\(s_1(X)\ge s_2(X)\ge\cdots\ge0\) denote its singular values, padded by zero.
A rectangular unitarily invariant seminorm is invariant under
\(X\mapsto UXV\) for unitary maps on the codomain and domain.  It is represented
by a symmetric gauge of the singular-value sequence.

The formal infrastructure needed by the paper consists of:
\begin{enumerate}
\item singular-value invariance under left and right unitary factors;
\item invariance under adjoint and under zero extension to a square operator;
\item the ideal inequalities
\[
  \Nn(CX)\le \norm{C}_{\op}\Nn(X),\qquad
  \Nn(XD)\le \Nn(X)\norm{D}_{\op};
\]
\item Ky Fan dominance: if
\(\sum_{j<k}s_j(X)\le\sum_{j<k}s_j(Y)\) for every \(k\), then
\(\Nn(X)\le\Nn(Y)\) for every unitarily invariant seminorm;
\item concrete operator, Frobenius, nuclear, Ky Fan, and Schatten instances.
\end{enumerate}

The 1970 proofs use the word ``compatible'' for norms satisfying the ideal
property with respect to the bound norm.  In finite dimensions, the required
compatibility follows from the singular-value ideal inequalities.

\section{Principal angles and the direct rotation}

Let \(P,Q\) be orthogonal projectors with ranges \(U,V\).  Their canonical
angles are encoded without choosing bases by the positive contractions obtained
from the two-projector decomposition.  In equal finite rank, the singular
values of \(QP\) are \(\cos\theta_j\), those of \(Q^\perp P\) are
\(\sin\theta_j\), and the nonzero singular values of \(P-Q\) are the sine
values with the standard multiplicities.

A complete Lean dictionary should include:
\begin{align*}
  \sv(QP)&=(\cos\theta_j)_j,\\
  \sv(Q^\perp P)&=(\sin\theta_j)_j,\\
  \norm{P-Q}_{\op}&=\max_j\sin\theta_j,\\
  \norm{P-Q}_{\HS}^2&=2\sum_j\sin^2\theta_j,\\
  \sv\bigl(2P^\perp QP\bigr)&=(\sin 2\theta_j)_j.
\end{align*}
Unequal ranks require explicit zero/one padding and a convention for unmatched
directions.

When the pair is acute, the direct rotation \(W\) is the canonical unitary
mapping \(U\) onto \(V\), acting as the identity on the common and doubly
orthogonal parts.  On each nontrivial principal two-plane it is ordinary
rotation by \(\theta_j\).  Its equivalent formulas include a polar-factor
formula, a trigonometric formula
\[
  W=\cos\Theta+J\sin\Theta,
\]
and the reflection identity
\[
  W^2=(2Q-I)(2P-I).
\]
The angle operator commutes with \(P,Q,J,W\).

\section{Compressions, Ritz operators, and residual identities}

Let \(X:F\to E\) be an isometric embedding representing an approximate
subspace.  For Hermitian \(A:E\to E\), define
\[
  M_X=X^*AX,\qquad R_X=AX-XM_X.
\]
More generally, for a proposed coordinate operator \(M\), define
\(R(A,X,M)=AX-XM\).  The following facts are implicit in the numerical-analysis
formulation and should be explicit Lean lemmas:
\begin{enumerate}
\item \(M_X\) is Hermitian;
\item \(X^*R_X=0\) (Galerkin orthogonality);
\item \(R_X=0\) iff \(\ran X\) reduces \(A\);
\item residuals transform naturally under a unitary change of coordinates;
\item if \(BX=XM\), then
\[
  R(A,X,M)=(A-B)X;
\]
\item in Frobenius norm,
\[
  \norm{R(A,X,M)}_{\HS}^2
  =\norm{R_X}_{\HS}^2+\norm{M_X-M}_{\HS}^2,
\]
so the Ritz compression minimizes the residual.
\end{enumerate}

Projecting the residual into an exact reducing decomposition gives the block
Sylvester equations from which all four angle estimates follow.

\section{The Sylvester estimate \texorpdfstring{$AX-XB=C$}{AX-XB=C}}

The elementary ordered theorem is the engine of the sharp constant-one
results.

\begin{theorem}[ordered Sylvester estimate]
Let \(A\) and \(B\) be Hermitian on possibly different finite-dimensional
Hilbert spaces.  Assume, after a common scalar translation, that
\(A\ge(c+\delta)I\) and \(B\le cI\), with \(\delta>0\).  If
\(AX-XB=C\), then for every compatible unitarily invariant seminorm,
\[
  \delta\,\Nn(X)\le\Nn(C).
\]
\end{theorem}

The proof first handles bound-norm coercivity and then passes to every
unitarily invariant norm via Ky Fan variational characterizations.  The roles
of \(A\) and \(B\) may be interchanged.  Interval/exterior separation is
reduced to ordered pieces and recombined through the projector/Ky Fan lemmas.

For arbitrary disjoint spectral sets separated by \(\delta\), the sharp
constant one generally fails.  The finite-dimensional Bhatia--Davis--McIntosh
extension instead has
\[
  \delta\,\Nn(X)\le \frac{\pi}{2}\Nn(C).
\]
This extension must remain a separate theorem so that it does not weaken the
classic interval theorem silently.

\section{The \texorpdfstring{$\sin\Theta$}{sin Theta} theorem}

Let \(U\) reduce \(A\), and let \(X:F\to E\) be an isometric approximate
basis with Hermitian coordinate operator \(M\).  Suppose
\(\spec(M)\subseteq[a,b]\), while
\(\spec(A|_{U^\perp})\) lies outside
\((a-\delta,b+\delta)\).  Projecting
\(AX-XM=R\) into \(U^\perp\) gives a Sylvester equation for
\(P_{U^\perp}X\).  Hence
\[
  \delta\,\Nn(P_{U^\perp}X)\le\Nn(R).
\]
The singular values of \(P_{U^\perp}X\) are the principal sines, producing the
residual \(\sin\Theta\) theorem.

If \(X\) spans a reducing subspace of \(B=A+H\), then
\(R=-HX\).  The ideal inequality and the square angle-operator dictionary give
the perturbation form
\[
  \delta\,\Nn(\sin\Theta)\le\Nn(H).
\]
The one-sided cross-projection and full projector-difference formulations must
be distinguished in unequal rank.  The general disjoint-spectrum variant has
the \(\pi/2\) constant described in the previous section.

\section{The \texorpdfstring{$\tan\Theta$}{tan Theta} theorem}

Assume an ordered gap between the unwanted exact block and the approximate
coordinate spectrum, and impose the Ritz/Galerkin condition so that the
relevant diagonal residual block vanishes.  The cosine block is then injective.
Writing the approximate subspace as the graph of an operator \(T:U\to U^\perp\),
the singular values of \(T\) are \(\tan\theta_j\).

The residual block equation and the ordered Sylvester estimate yield
\[
  \delta\,\Nn(\tan\Theta)\le\Nn(R)
\]
for every unitarily invariant norm.  In the perturbation formulation, the
vanishing compression condition is
\[
  QHQ=0
\]
on the chosen perturbed subspace, and the conclusion is
\[
  \delta\,\Nn(\tan\Theta)\le\Nn(H).
\]
The proof needs explicit lemmas establishing transversality, the graph
representation, the tangent singular-value dictionary, and the Ky Fan passage
from finite-rank test projectors to arbitrary unitarily invariant norms.

\section{The \texorpdfstring{$\sin 2\Theta$}{sin 2 Theta} theorem}

The double-angle theorem depends only on the internal gap of the unperturbed
operator.  Let \(U\) reduce \(A\), and assume the spectra of
\(A|_U\) and \(A|_{U^\perp}\) are separated by \(\delta\).  For an approximate
invariant pair \((X,M)\), the two projected residual equations combine to give
\[
  \delta\,\Nn(\sin 2\Theta)\le 2\Nn(R).
\]

For the perturbation form, let \(V\) reduce \(B=A+H\).  Reflect \(A\) through
\(V\):
\[
  A'=J_VAJ_V,\qquad D_V(A)=A'-A.
\]
The cross-projection between \(U\) and its reflected image has singular values
\(\sin 2\theta_j\).  The single-angle theorem gives
\[
  \delta\,\Nn(\sin 2\Theta)\le\Nn(D_V(A)).
\]
Since \(J_VBJ_V=B\),
\[
  D_V(A)=J_V(A-B)J_V-(A-B),
\]
and unitary invariance plus the triangle inequality gives
\(\Nn(D_V(A))\le2\Nn(H)\).  Therefore
\[
  \delta\,\Nn(\sin 2\Theta)\le2\Nn(H).
\]

\section{The \texorpdfstring{$\tan 2\Theta$}{tan 2 Theta} theorem}

Assume in addition that the perturbation is off diagonal relative to the
unperturbed splitting:
\[
  PHP=0,\qquad P^\perp H P^\perp=0.
\]
The coupled block equations become a Riccati/graph equation.  On each principal
two-plane the denominator contains \(\cos 2\theta_j\), and the off-diagonal gap
forces the selected solution to satisfy \(\theta_j<\pi/4\).  The resulting
sharp estimates are
\[
  \delta\,\Nn(\tan 2\Theta)\le2\Nn(R),
  \qquad
  \delta\,\Nn(\tan 2\Theta)\le2\Nn(H).
\]
A complete formal proof requires more than the largest-angle scalar algebra:
it needs the tangent-double-angle singular-value dictionary, Ky Fan estimates,
and a theorem selecting the acute reducing subspace of the perturbed operator.

\section{Spectral selection and repulsion}

The interpretation section of the 1970 paper proves that an off-diagonal
perturbation admits a unique ``correct'' reducing subspace on the prescribed
side of the gap.  In finite dimensions the formal theorem should construct this
subspace as a canonical spectral subspace of \(A+H\), prove that it is acute
relative to \(U\), and prove the equivalence between the spectral-side
conditions and the bound \(\Theta\le\pi/4\).

The same analysis yields spectral repulsion.  The perturbed eigenvalues on each
side of the gap move away from the gap, and compression through the cosine
block gives both ordered-eigenvalue and symmetric-gauge inequalities.  The Lean
roadmap therefore includes:
\begin{enumerate}
\item existence and uniqueness of the reducing subspace preserving the gap;
\item acuteness and strict \(\theta_{\max}<\pi/4\) under the quantitative
hypotheses;
\item compression comparison through the cosine block;
\item eigenvalue-by-eigenvalue repulsion;
\item a Ky Fan/UI-norm repulsion statement;
\item the continuation argument that keeps the canonical spectral projector on
the same branch along \(A+tH\).
\end{enumerate}

\section{Extremal properties of the direct rotation}

Among all unitaries carrying \(U\) onto \(V\), the direct rotation has the
precise extremal properties isolated in Section~4.  The unconditional
all-unitarily-invariant-norm statement is for the positive displacement
square \((I-W^*)(I-W)\), not for \(I-W\) itself.  The latter minimization can
fail for some unitarily invariant norms; Davis--Kahan recover it uniformly
when the largest angle is at most \(\pi/3\).  The finite roadmap records:
\begin{enumerate}
\item uniqueness from the intertwining relation plus nonnegative real part;
\item minimization of \(\Nn{(I-W^*)(I-W)}\) for every unitarily invariant norm;
\item minimization of \(\Nn{I-W}\) for every unitarily invariant norm under
      the largest-angle bound \(\theta_{\max}\le\pi/3\);
\item minimization of maximal vector displacement;
\item minimization of the sum of squared displacements of an orthonormal basis;
\item compatibility with the blockwise polar-factor intertwining unitary used
in Davis's 1963 total-rotation theory.
\end{enumerate}
These results are geometric rather than prerequisites for the four inequalities,
but they complete the classic theory's interpretation of the angle operator.

\section{Sharpness and planar equality models}

All four constants are best possible.  The equality calculations reduce to a
real or complex two-dimensional principal plane.  Take
\[
  P=\begin{bmatrix}1&0\\0&0\end{bmatrix},\qquad
  Q=\begin{bmatrix}\cos^2\theta&\sin\theta\cos\theta\\
                    \sin\theta\cos\theta&\sin^2\theta\end{bmatrix},
\]
and diagonal or off-diagonal Hermitian operators chosen so that the gap and
residual equations are equalities.  Direct sums of these planes realize any
finite singular-value list and therefore equality simultaneously for all
unitarily invariant norms.

The Lean sharpness layer should provide the explicit matrices/operators,
compute their principal angles and residuals, prove equality in each theorem,
and derive that the constants one and two cannot be decreased.  It should also
record the pole obstruction for \(\tan\Theta\) and the \(\pi/4\) branch
restriction for \(\tan2\Theta\).

\section{Lean module map and proof order}

\begin{longtable}{p{0.36\textwidth}p{0.56\textwidth}}
\toprule
Lean module & Mathematical responsibility \\
\midrule
\texttt{Basic.lean} & canonical spectral subspaces, projectors, angle objects,
reducing/gap predicates, pinch and off-diagonal predicates \\
\texttt{RectangularUINorm.lean} & rectangular UI seminorms, singular-value
instances, ideal property, Fan dominance and square compatibility \\
\texttt{Residual.lean} & compression, Ritz operator, residual algebra,
Galerkin orthogonality, graph/sine/tangent embeddings \\
\texttt{Sylvester.lean} & ordered and interval/exterior \(AX-XB=C\) estimates,
Ky Fan form, and the separate \(\pi/2\) extension \\
\texttt{SinTheta.lean} & residual, perturbation, projector, interval, Ky Fan,
operator, and Frobenius \(\sin\Theta\) statements \\
\texttt{TanTheta.lean} & transversality, graph operator, tangent singular values,
residual and perturbation \(\tan\Theta\) statements \\
\texttt{SinTwoTheta.lean} & mirror defect and residual/perturbation
\(\sin2\Theta\) statements \\
\texttt{TanTwoTheta.lean} & off-diagonal graph/Riccati theory,
\(\tan2\Theta\), spectral selection, acuteness, and repulsion \\
\texttt{Generalized.lean} & Theorems 6.1--6.3 and 8.2:
non-orthonormal trial maps, unequal dimensions, square-norm fallback, and
continuation to the acute branch \\
\texttt{DirectRotation.lean} & canonical direct rotation and extremal geometry \\
\texttt{Davis1963.lean} & total rotation, eigenvalue-motion subtraction, and
one-vector double-angle ancestors \\
\texttt{Sharpness.lean} & two-dimensional extremizers and optimal constants \\
\texttt{Statistics.lean} & population-gap, aligned-basis, and
Yu--Wang--Samworth corollaries \\
\texttt{SingularSubspace.lean} & Gram operators, Hermitian dilation, and
left/right singular-subspace corollaries \\
\bottomrule
\end{longtable}

The intended implementation order is the order of this table.  Later headline
theorems should be proved by the named lower-level seams rather than by
repeating eigenbasis calculations.  Existing proved files should be reused
where their statement matches; scaffold declarations may then become wrappers
or be deleted in favor of the proved canonical theorem.

\begin{thebibliography}{9}
\bibitem{Davis1963}
C. Davis,
\emph{The rotation of eigenvectors by a perturbation},
J. Math. Anal. Appl. 6 (1963), 159--173.

\bibitem{DavisKahan1970}
C. Davis and W. M. Kahan,
\emph{The rotation of eigenvectors by a perturbation. III},
SIAM J. Numer. Anal. 7 (1970), 1--46.

\bibitem{Mirsky1960}
L. Mirsky,
\emph{Symmetric gauge functions and unitarily invariant norms},
Quart. J. Math. Oxford 11 (1960), 50--59.

\bibitem{YWS2015}
Y. Yu, T. Wang, and R. J. Samworth,
\emph{A useful variant of the Davis--Kahan theorem for statisticians},
Biometrika 102 (2015), 315--323.
\end{thebibliography}

\end{document}
